\documentclass[letterpaper,twocolumn,10pt]{article}
\usepackage{garb}

\renewcommand\abstractname{\textbf{Abstract}}

%-------------------------------------------------------------------------------
\begin{document}

%don't want date printed
\date{}

% make title bold and 14 pt font (Latex default is non-bold, 16 pt)
\title{GARB Jonathan + Elliot}

\author{{\rm Elev Elevkvist} \\
\textit{Teknikprogrammet} \\
\textit{Maja Beskowgymnasiet}\\
{\rm Umeå, Sverige}}

\maketitle

\begin{abstract}
The abstract is written in English. Focus on your improvement study: what
product feature (X) you measured, what method (Y) you used to improve it, and
what your specific product (Z) was. Include your key findings - was method Y
effective for improving X on your product type Z?

Your abstract should:
\begin{itemize}
    \item consist of a single paragraph of 150-250 words, with correct grammar
    and unambiguous terminology;
    \item be self-contained with no abbreviations, footnotes, references, or
    mathematical equations;
    \item highlight the practical implications of your improvement study;
    \item mention your research question, measurement methodology, key results,
    and conclusions about the effectiveness of your improvement method.
\end{itemize}
\end{abstract}

\selectlanguage{swedish}

\section{Inledning}

Laddningstider utgör en central del av användarupplevelsen i moderna
mobilapplikationer. Tidigare forskning visar att användare är mycket känsliga
för även små fördröjningar; svarstider som överstiger en sekund bryter
användarens tankegång, medan fördröjningar på tio sekunder eller mer markant
ökar sannolikheten för att applikationen överges \cite{nielsen1993response}. För
datadrivna applikationer, där innehåll kontinuerligt hämtas från externa
databaser, utgör nätverksanrop en betydande del av den totala laddningstiden.

Trots denna betydelse saknas tydliga riktlinjer för hur databasanrop bör
struktureras i Flutter- och Firebase-baserade applikationer för att minimera
latens. Detta skapar ett behov av att empiriskt undersöka hur olika
implementationsstrategier påverkar laddningstider i praktiken. Inom
systemdesign och operativsystem är det väl etablerat att effektiv hantering av
I/O-operationer är avgörande för systemets prestanda
\cite{arpachiDusseau18:osbook}. %Läs igenom och ta möjligtvis bort

%%% Bör ändra den här texten så att den endast beskriver egenskapen X, men
%bibehålla strukturen + första källan. Kan även vara bra att kolla hur abstaktet
%är formulerat.

Beskriv det problem eller behov som motiverar din förbättring. Referera till
vetenskapligt publicerad litteratur inom området\cite{arpachiDusseau18:osbook}
och tidigare forskning om liknande förbättringsmetoder.

\subsection{Produktbeskrivning}

Produkten som studeras är mobilapplikationen \textit{Neatify}, utvecklad med
ramverket Flutter och med Firebase Firestore som databashanteringssystem.
Flutter möjliggör plattformsoberoende utveckling där applikationslogik skrivs i
programmeringsspråket Dart \cite{TODO}.

Neatify syftar till att förbättra användarens produktivitet genom
spelifiering, där användaren motiveras att utveckla positiva beteenden genom
belöningar och kontinuerlig feedback baserad på interaktioner i applikationen.

Applikationen består av flera vyer som dynamiskt hämtar användardata från
Firestore för att rendera innehåll och funktionalitet. Vid navigering mellan
dessa vyer utförs flera oberoende nätverksanrop. Särskilt sidorna
\textit{Profile} och \textit{Leaderboard} innehåller ett större antal sådana
anrop och utgör därför fokus för denna studie.

\subsection{Förbättringsmetod}
\label{sec:forbattringsmetod}

Den metod som undersöks för att förbättra laddningstiderna i Neatify är
parallellisering av nätverksanrop med hjälp av \texttt{Future.wait()} i Dart.
Istället för att exekvera oberoende anrop sekventiellt, där varje anrop väntar
på att det föregående ska slutföras, initieras samtliga anrop samtidigt och
applikationen inväntar deras gemensamma resultat innan rendering sker
\cite{dart_future_wait}.

Metoden bygger på principer inom asynkron programmering, där oberoende
I/O-operationer kan initieras parallellt för att minska den totala väntetiden
\cite{dart_async}. Detta angreppssätt är teoretiskt motiverat då den totala
exekveringstiden i nätverksberoende system ofta begränsas av den långsammaste
operationen snarare än summan av alla operationer, vilket gör parallellisering
till en naturlig strategi för att reducera latens. 

\subsection{Syfte och frågeställning}

Syftet med denna rapport är att empiriskt utvärdera effekten av parallellisering av nätverksanrop i Flutter- och Firebase-baserade mobilapplikationer, med fokus på laddningstider i applikationen Neatify. Utvärderingen avgränsas till sidorna \textit{Profile} och \textit{Leaderboard}, där laddningstid mäts med median och p95.

\textbf{Huvudfrågeställning:}
\begin{itemize}
    \item I vilken utsträckning kan parallellisering av nätverksanrop förbättra laddningstider i Flutter- och Firebase-baserade mobilapplikationer?
\end{itemize}

\textbf{Stödjande frågeställningar:}
\begin{itemize}
    \item Hur stor förbättring av laddningstider (median och p95) uppnås efter implementering av parallellisering av nätverksanrop på sidorna \textit{Profile} och \textit{Leaderboard}?
    \item Hur påverkas resultatet av databasstorlek och den nätverkslatens som observerades under mätningarna?
    \item Är förbättringen statistiskt signifikant och praktiskt relevant?
\end{itemize}

\subsection{Teori}

För att förstå parallelliseringens potentiella effekt på laddningstider är det viktigt att först förstå de bakomliggande faktorerna:
\begin{itemize}
    \item \textbf{Flutter}: Flutter är Googles gratis och open-source UI-ramverk för att bygga nativt kompilerade applikationer för mobil, webb och desktop från en enda kodbas. \cite{flutter_faq}
    \item \textbf{Firebase}: Firebase är en plattform för mobil- och webbapplikationsutveckling som erbjuder en rad tjänster, inklusive realtidsdatabaser, autentisering och hosting. \cite{firebase_docs} De specifika tjänster som används i Neatify är Firestore, en NoSQL-databas som möjliggör skalbar och flexibel datalagring, och Firebase Authentication för att hantera användarautentisering. \cite{firestore_docs} \cite{firebase_auth_docs} Men även Firebase Performance Monitoring, som används för att mäta och analysera applikationens prestanda i realtid. \cite{firebase_perf_mon_docs}
\end{itemize}

%Kolla igenom det finns mer teori som är relevant att ta med, t.ex. om asynkron programmering i Dart, eller om nätverkslatens och dess påverkan på användarupplevelse.

\subsection{Bakgrund}

Parallellisering av databasanrop har sin teoretiska grund i forskningsområdet
parallel query processing, där målet är att minska exekveringstid genom att
utföra flera operationer samtidigt. Tidig forskning inom området visar att
databassökningar kan delas upp i deluppgifter som utförs parallellt över flera
processorer, vilket möjliggör prestandavinster jämfört med sekventiell
exekvering \cite{dewitt1992parallel}.

I moderna databassystem har dessa principer vidareutvecklats i distribuerade
miljöer, där data partitioneras över flera noder och bearbetas parallellt innan
resultaten sammanställs. Studier visar att distribuerad parallell bearbetning kan
förbättra skalbarhet och svarstider, särskilt när partitioneringen utformas för
att minimera kommunikation mellan noder, exempelvis genom grafbaserade metoder.
Samtidigt påverkas vinsterna av faktorer som kommunikationskostnader och synkronisering
mellan noder \cite{nam2019parallel}.

Forskning inom samtidiga databassystem, inklusive GPU-baserade lösningar, visar
att flera förfrågningar kan exekveras parallellt för att bättre utnyttja
tillgängliga beräkningsresurser, även om detta kräver effektiv hantering av
resurstilldelning och exekveringsplanering \cite{li2019concurrent}. Vidare
indikerar studier att parallella exekveringstekniker kan bidra till förbättrad
skalbarhet genom att dela upp operationer i samtidiga deluppgifter
\cite{zhu2025research}.
%Måste överväga om den här källan är en hjälp eller om den skadar rapportens trovärdighet.

Även om dessa studier främst behandlar CPU-bundna eller distribuerade databassystem, bygger de på grundläggande principer om parallellism, väntetider och resursutnyttjande. Dessa principer är relevanta även i mobilapplikationer, där databasanrop ofta utgörs av nätverksbaserade I/O-operationer. I sådana fall påverkas svarstiden av faktorer som nätverkslatens och serverrespons, och genom att initiera flera oberoende anrop parallellt kan den totala väntetiden potentiellt reduceras. Därmed kan mobilapplikationers databasanrop analyseras utifrån samma övergripande prestandaprinciper som beskrivs i tidigare forskning.
%BÖR STÄRKAS MED KÄLLA FÖR MOBILA APPLIKATIONER eller MOBILER

\section{Metod}

\subsection{Implementering av förbättring}

Implementeringen av förbättringen bestod av att ersätta de sekventiella 
Firestore-anropen på \textit{Profile page} och \textit{Leaderboard page} med 
parallella anrop med hjälp av Darts \texttt{Future.wait()}, vars funktion 
beskrivs närmare i avsnitt \ref{sec:forbattringsmetod}.

Förändringen genomfördes uteslutande i de delar av koden som hanterar 
datahämtning vid sidladdning, utan ändringar i UI-logik, datamodeller eller 
övrig applikationslogik, för att säkerställa att eventuella skillnader i 
mätresultaten kan tillskrivas enbart den parallella exekveringen. Implementationen 
versionshanterades i branch \texttt{parallel-db}, isolerad från den ursprungliga 
koden i \texttt{staging}.

Innan de fullskaliga utvärderingsmätningarna genomfördes verifierades 
implementeringen manuellt genom att navigera till de berörda sidorna och 
bekräfta att all data laddades korrekt och att sidorna renderades som förväntat 
efter övergången till parallell exekvering.

\subsection{Genomförande av mätningar}

För att kunna besvara frågeställningen togs ett test-script, som använder
Flutters \texttt{IntegrationTestWidgetsFlutterBinding} med en initial cold-boot
för varje sida där firestores caching är inaktiverat, fram för att snabbt och
enkelt kunna köra flera mätningar av laddningstider, som mäts från att sidan
öppnas till att allt innehåll på sidan är renderat och interaktivt. I
kombination med detta test-script kopierades de två relevanta sidorna,
\textit{Profile page} och \textit{Leaderboard page}, till samma branch som den
parallella implementeringen för att kunna köra ett alternerande mönster av
mätningar, där jämna mätningspar kördes med den sekventiella versionen först och
udda mätningspar kördes med den parallella versionen först, för att minimera
påverkan av nätverksvariation på resultaten. Med test-scriptet gjordes 30
mätningar av vardera sida för både den sekventiella och parallella versionen,
vilket ger en total av 120 mätta warm-laddningstider.

För att säkerställa att mätningarna var så stabila som möjligt, så kördes alla
mätningar i en kontrollerad miljö där inga andra program var igång och där
telefonen inte användes under mätningarna. Tiderna från mätningarna loggades i
en csv-fil för att sedan kunna analyseras med hjälp av Python och Matplotlib. 

Mätningarna upprepades med olika storlekar på databasen för att kunna
utvärdera förbättringens effekt på olika belastningsnivåer. För att kunna
bestämma databasens storlek användes ett seeding-script, skrivet i JavaScript
med Node.js, som genererade testdata för databasen i storlekarna 100, 500, 1000
användare. För att säkerställa att databasen var samma, förutom för
användarantalet, sattes alla andra fields till samma värde för alla mätningar.
För varje databasstorlek antecknades både storleken och nätverkshastigheten för
att kunna kontrollera dessa faktorer i analysen.

För alla mätningar användes en OnePlus 12R med Android 16 och OxygenOS 16.0.3
som testenhet. Testenheten var ansluten till en stabil Wi-Fi-anslutning och hade
alla onödiga bakgrundsprocesser stängda för att minimera störningar. För att
säkerställa att resultaten var så tillförlitliga som möjligt, så genomfördes
alla mätningar under liknande förhållanden, inklusive tid på dagen och
batterinivå. 

För att säkerställa reproducerbarhet och transparens dokumenterades den exakta
kodstatusen som användes under de initiala mätningarna via Git. Branch
\texttt{parallell-db} innehåller den parallella implementationen tillsammans med
kopior av sidorna med sekventiella anrop. Men den innehåller även seeding- och
test-scriptet. Dock kräver seeding-scriptet ett service account med behörighet
att skriva till Firestore, vilket av säkerhetsskäl inte kan inkluderas i det
publika Git-repot. Av denna anledning behöver en reproducerare av studien skapa
ett egnet firebase-projekt, generera en egen serviceAccountKey.json och köra
seeding-scriptet för att återskapa en funktionellt identisk testmiljö.
Instruktioner för detta finns dokumenterade i repots README på branch
\texttt{parallel-db}.

\subsection{Statistisk analys}

Den statistiska analysen utgick från att varje mätning av den sekventiella implementationen matchades mot motsvarande mätning av den parallella implementationen för samma sida, iteration och databasstorlek. Analysen genomfördes därför som en parad jämförelse, där skillnaden i laddtid definierades som sekventiell minus parallell laddtid. Positiva differenser indikerar därmed kortare laddtid för parallella anrop.

Inför analysen exkluderades cold-boot-mätningar från den inferentiella
jämförelsen, eftersom dessa i hög grad påverkas av initialisering och
cache-relaterade uppstartseffekter och därmed inte speglar den stabila
laddningsprestandan vid återkommande sidvisningar. Den primära analysen
baserades i stället på warm-mätningar. Saknade värden hanterades med pairwise
deletion, vilket innebär att endast kompletta par inkluderades i
hypotesprövningen. Normalisering användes inte, eftersom den parade designen
redan hanterar variationer mellan mätningar.

Som beskrivande statistik rapporterades antal mätningar, medelvärde, median och spridningsmått för respektive sida och implementation. För hypotesprövning användes ett parat Wilcoxon signed-rank-test, valt eftersom mätningarna är beroende och laddtider kan vara snedfördelade samt innehålla avvikande mätningar. Testet genomfördes separat för \textit{Profile page} och \textit{Leaderboard page} för att undvika att sidspecifika beteenden sammanblandas i ett gemensamt test. Signifikansnivån sattes till $\alpha = 0.05$ och den statistiska prövningen kompletterades med en bedömning av praktisk signifikans baserad på observerade skillnader i millisekunder.

\section{Resultat}

\subsection{Mätningar av sekventiella anrop}

Presentera resultaten från dina initiala mätningar av egenskap X på den
ursprungliga produkten och referera till dem i Tabell \ref{tab:baseline}

Utifrån mätdatan för de sekventiella anropen, innan implementeringen av
förbättringen, samanställdes tabellen nedan som visar beskrivande statistik för
varje databasstorlek och sida. Tabellen inkluderar medelvärde,
standardavvikelse, 95\% konfidensintervall, nätverkslatens och p95-värden för
både profil- och leaderboard-sidorna.
%%%Dessa data ger en grundläggande förståelse för prestandan hos de sekventiella anropen innan förbättringen implementerades.

\begin{table}[H]
  \centering
  \small
  \caption{Beskrivande statistik för sekventiella anrop, uppdelat på databasstorlek och sida}
    \begin{tabular}{lccccccc}
        \toprule
          Databasstorlek (users)     & Sida & n & Medelvärde (ms) &
          Standardavvikelse (ms) & 95\% KI (ms) & Nätverkslatens (ms), medel ±
          std & p95 (ms)    \\
        \midrule
            100         & Profile page & 150 & 699.98 & 66.61 &
            689.23--710.73 & 22.00 $\pm$ 3.16 & 807.10 \\
            100         & Leaderboard page  & 150 & 570.88 & 33.69 &
            565.44--576.32 & 22.00 $\pm$ 3.16 & 612.00 \\
            500         & Profile page & 150 & 703.55 & 65.67 &
            692.96--714.15 & 22.80 $\pm$ 3.27 & 807.55  \\
            500         & Leaderboard page  & 150 & 563.13 & 48.28 &
            555.34--570.92 & 22.80 $\pm$ 3.27 & 626.55 \\
            1000         & Profile page & 150 & 683.95 & 54.47 &
            675.16--692.74 & 23.80 $\pm$ 3.63 & 797.65 \\
            1000         & Leaderboard page  & 150 & 567.45 & 39.21 &
            561.12--573.77 & 23.80 $\pm$ 3.63 & 613.20 \\
        \bottomrule
    \end{tabular}
  \label{tab:sekventiell}
\end{table}

\subsection{Mätningar av parallella anrop}

På samma sätt som för de sekventiella anropen samanställdes tabellen nedan för
resultaten för de parallella anropen, efter implementering av förbättringen.

\begin{table}[H]
  \centering
  \small
  \caption{Beskrivande statistik för parallella anrop, uppdelat på databasstorlek och sida}
    \begin{tabular}{lccccccc}
        \toprule
          Databasstorlek (användare)     & Sida & n & Medelvärde (ms) &
          Standardavvikelse (ms) & 95\% KI (ms) & Nätverkslatens (ms), medel ±
          std & p95 (ms)    \\
        \midrule
            100         & Profile page & 150 & 480.63 & 56.67 &
            471.49--489.78 & 22.00 $\pm$ 3.16 & 580.10 \\
            100         & Leaderboard page  & 150 & 459.33 & 35.98 &
            453.52--465.13 & 22.00 $\pm$ 3.16 & 521.90 \\
            500         & Profile page & 150 & 495.91 & 68.91 &
            484.79--507.03 & 22.80 $\pm$ 3.27 & 596.55  \\
            500         & Leaderboard page  & 150 & 460.19 & 52.03 &
            451.79--468.58 & 22.80 $\pm$ 3.27 & 584.40 \\
            1000         & Profile page & 150 & 473.54 & 60.17 &
            463.83--483.25 & 23.80 $\pm$ 3.63 & 583.10 \\
            1000         & Leaderboard page  & 150 & 459.88 & 43.05 &
            452.93--466.83 & 23.80 $\pm$ 3.63 & 545.15 \\
        \bottomrule
    \end{tabular}
  \label{tab:parallel}
\end{table}

\subsection{Jämförande analys}

\textbf{Förbättringsresultat:}

Med värdena från Tabell \ref{tab:sekventiell} och Tabell \ref{tab:parallel}
kunde förbättringen, absolut och relativ, beräknas för varje databasstorlek och
sida. Effektstorleken kunde också beräknas, vilket gjordes med rangbiseriell
korrelation (r_rb),
%Detta bör gå till metod: vilket är lämpligt för icke-parametriska data och ger en intuitiv tolkning av förbättringens styrka.
och de specifika resultaten för varje databasstorlek och sida redovisas nedan:

\begin{table}[H]
  \centering
  \small
  \caption{Jämförelse före och efter}
    \begin{tabular}{lcccc}
        \toprule
            Databasstorlek (användare)  & sida  & absolut förbättring (ms)    & relativ förbättring & r_rb   \\
        \midrule
            100 & Profile page  & 219.35   & 31.34\% & 1.000   \\
            100 & Leaderboard page    & 111.55    & 19.54\% & 0.995   \\
            500 & Profile page  & 207.64   & 29.51\% & 1.000   \\
            500 & Leaderboard page    & 102.95    & 18.28\% & 0.960   \\
            1000 & Profile page  & 210.41   & 30.76\% & 1.000   \\
            1000 & Leaderboard page    & 107.57    & 18.96\% & 0.965   \\
        \bottomrule
    \end{tabular}
  \label{tab:comparison}
\end{table}

\subsection{Statistisk signifikans}

Med den statiska analysen med Wilcoxon signed-rank test 
%%%Koppla till metod + utöka möjligen
kunde p-värden bestämmas för de olika databasstorlekarna och sidorna.

\begin{table}[H]
  \centering
  \small
  \caption{Jämförelse före och efter}
    \begin{tabular}{lccc}
        \toprule
            Databasstorlek (användare)  & sida  & W-statistik    & p-värde   \\
        \midrule
            100 & Profile page  & 11325.00   & 1.14672 $\times 10^{-26}$  \\
            100 & Leaderboard page    & 11294.00    & 2.12553 $\times 10^{-26}$   \\
            500 & Profile page  & 11325.00   & 1.1469 $\times 10^{-26}$   \\
            500 & Leaderboard page    & 10949.50    & 1.46696 $\times 10^{-24}$   \\
            1000 & Profile page  & 11325.00   & 1.14655 $\times 10^{-26}$   \\
            1000 & Leaderboard page    & 11125.50    & 5.89925 $\times 10^{-25}$   \\
        \bottomrule
    \end{tabular}
  \label{tab:wilcoxon}
\end{table}

Det resultat som presenteras i Tabell \ref{tab:wilcoxon} visar att alla
förbättringar är statistiskt signifikanta, med p-värden under den konventionella
tröskeln på 0.05. 

\subsection{Implementeringsanalys}

Under implementeringen av parallella nätverskanrop gjordes följande
observationer:
\begin{itemize}
    \item Implementeringstid: uppskattningsvis under 1 timme (ingen exakt
    tidlogg fördes)
    \item Kodkomplexitet: Profil-sidan: +2 rader kod, Leaderboard-sidan: +4
    rader kod
    \item Bieffekter: Inga tydliga bieffekter observerades under testning
    \item Kompatibilitet: Testad i Flutter + Firebase-miljö utan
    funktionsbortfall.
\end{itemize}

\section{Diskussion}

I denna del ska du diskutera och utvärdera dina resultat i relation till syftet
och frågeställningarna.



\subsection{Utvärdering av förbättringens effektivitet}

De uppmätta förbättringarna i Tabell \ref{tab:comparison} visar att den absoluta
förbättringen i laddningstid var mellan 102.95 ms och 219.35 ms, vilket
motsvarar en relativ förbättring på mellan 18.28\% och 31.34\% beroende på
databasstorlek och sida. Dessa förbättringar är inte bara statistiskt
signifikanta, som visas i Tabell \ref{tab:wilcoxon}, utan också praktiskt
betydelsefulla, särskilt i kontexten av mobilapplikationer där användare är
känsliga för laddningstider. Även om laddtiderna höll sig under den kritiska
gränsen på 1 sekund \cite{nielsen1993response}, kan en förbättring på över 100
ms fortfarande ha en positiv inverkan på användarupplevelsen, särskilt när det
gäller att minska känslan av latens och förbättra den övergripande
responsiviteten i applikationen. 

\subsection{Effekt av databasstorlek och nätverkslatens}

Förbättringen var konsekvent över samtliga testade databasstorlekar, vilket
tyder på att parallellisering av nätverksanrop är en robust metod för att
förbättra laddningstider vid olika belastningsnivåer. En svag tendens till
större absolut förbättring observerades för den minsta databasen, men utifrån
den insamlade datan kan detta inte tolkas som en tydlig trend.

Nätverkslatensen, som var relativt stabil över mätningarna, verkar inte ha haft
någon större inverkan på förbättringens storlek. Dock kan inte någon slutsats
dras om interaktionen mellan nätverkslatens och förbättringens effektivitet
baserat på den data som samlats in, eftersom nätverkslatensen inte varierade
tillräckligt mycket för att möjliggöra en sådan analys. 

\textbf{Förbättringens storlek:}
Den uppmätta förbättringen på 23.6\% kan betraktas som substantiell inom detta
område. Jämfört med tidigare studier på liknande produkter\cite{waldspurger02}
ligger detta resultat i linje med förväntningar för denna typ av optimering.

\textbf{Statistisk tillförlitlighet:}
Med p < 0.001 och en effektstorlek (Cohen's d) på 2.74 visar resultaten på en
stark och tillförlitlig effekt. Det stora antalet mätningar (n=35) och den låga
standardavvikelsen stärker tillförlitligheten i resultaten.

\subsection{Lämplighet av parallella nätverksanrop för en Flutter- och Firebase-baserad applikation}

Parallellisering av nätverksanrop gav tydliga prestandaförbättringar i form av
kortare laddningstider och därmed en bättre användarupplevelse i Neatify.
Metoden krävde endast begränsade kodändringar och påverkade inte applikationens
funktionalitet eller stabilitet. Samtidigt förutsättet metoden att
nätverksanropen är oberoende, och den kan göra felhantering något mer komplex,
vilket är en viktig aspekt att överväga i mer komplexa applikationer än Neatify. 

\textbf{Fördelar:}
\begin{itemize}
    \item Betydligt förbättrade laddningstider, vilket kan förbättra användarupplevelsen
    \item Relativt enkel implementering med begränsad kodomstrukturering
    \item Kompatibel med befintlig arkitektur
    \item Låg risk för introduktion av nya buggar
\end{itemize}

\textbf{Nackdelar och begränsningar:}
\begin{itemize}
    \item Ökad minnesanvändning kan vara problematisk för resursbegränsade
    enheter
    \item Implementering kräver fördjupad kunskap inom området
    \item Underhållskostnader kan öka på grund av ökad komplexitet
\end{itemize}

\subsection{Metodkritik och begränsningar}

Metoden bedöms ha varit lämplig för att besvara frågeställningen, eftersom den
visade en statistiskt signifikant förbättring i laddningstider efter
implementeringen av parallella nätverksanrop i den testade applikationen.
Samtidigt finns det begränsningar som bör beaktas vid tolkningen av resultaten.
Mätningarna genomfördes i en kontrollerad miljö med en specifik enhet, stabil
Wi-Fi-anslutning och inaktiverad cachelagring. Detta ger en tydlig bild av
förbättringens effekt under kontrollerade förhållanden, men innebär också att
resultaten inte fullt ut speglar den variation som kan förekomma i verkliga
användningsförhållanden, där enheter, nätverksförhållanden och användarbeteenden
kan variera kraftigt.

Eftersom resultaten är knutna till Neatify och den aktuella implementationen av
parallella nätverksanrop bör slutsatser om generaliserbarhet dras med
försiktighet. Andra applikationer, särskilt de med mer komplexa beroenden mellan
nätverksanrop eller som körs på mer resursbegränsade enheter, kan ge andra
resultat.

\textbf{Testmiljöns representativitet:}
Mätningarna genomfördes i en kontrollerad miljö som kanske inte helt återspeglar
verkliga användningsförhållanden.

\textbf{Generaliserbarhet:}
Resultaten är specifika för den testade produkttypen och implementeringen.
Generalisering till andra produkttyper bör göras med försiktighet.

Då testmiljön endast bestod av en testenhet och en stabil Wi-Fi-anslutning kan
resultaten vara mindre generaliserbara till en bredare användarbas med
varierande enheter och nätverksförhållanden. 

\subsection{Praktiska implikationer}

Jämför dina resultat med tidigare forskning eller liknande arbeten och notera
skillnader och likheter.

Resultaten visar att parallellisering av oberoende Firestore-anrop kan ge
tydliga och praktiskt relevanta förbättringar i laddningstid för Flutter- och
Firebase-baserade applikationer. För Neatify innebär detta att sidan kan
renderas snabbare utan att applikationens funktionalitet behöver förändras i
grunden. Metoden är därför särskilt lämplig i situationer där flera databasanrop
kan utföras oberoende av varandra och där snabbare sidladdning är viktig för
användarupplevelsen.

Resultaten ligger i linje med tidigare forskning om parallell exekvering och
I/O-bundna operationer, där total väntetid kan minskas genom att oberoende
arbete utförs samtidigt i stället för sekventiellt. Samtidigt visar studien att
förbättringens storlek varierar mellan olika sidor och databasstorlekar, vilket
innebär att parallellisering inte bör betraktas som en universallösning, utan
som en metod som bör användas när beroendena mellan anropen tillåter det.

\textbf{Rekommendationer för implementation:}
Baserat på studiens resultat rekommenderas parallellisering av oberoende
nätverksanrop i Flutter- och Firebase-baserade applikationer, särskilt när
kortare laddningstid prioriteras och när anropen inte är beroende av varandras
resultat.

\textbf{Förslag på framtida forskning:}
\begin{itemize}
\item Utvärdera metoden på fler typer av Flutter-applikationer med annan datastruktur och andra beroenden mellan anrop.
\item Undersöka hur resultaten påverkas av varierande nätverksförhållanden och olika hårdvarukonfigurationer.
\item Jämföra parallellisering med andra optimeringsmetoder, till exempel cachning eller reducering av antal anrop.
\item Genomföra användarstudier för att undersöka hur de uppmätta förbättringarna upplevs i praktiken.
\end{itemize}

\subsection{Slutsats}

Studien visar att parallellisering av oberoende nätverksanrop med Future.wait()
är en effektiv metod för att förbättra laddningstider i den testade Flutter- och
Firebase-applikationen Neatify. Jämfört med sekventiell exekvering gav metoden
genomgående kortare laddningstider för både Profile page och Leaderboard page,
med absoluta förbättringar på cirka 103-219 ms och relativa förbättringar på
cirka 18-31 procent beroende på sida och databasstorlek.

De observerade skillnaderna var statistiskt signifikanta i samtliga jämförelser
enligt Wilcoxon signed-rank-test. Resultaten var dessutom stabila över de
testade databasstorlekarna, utan någon tydlig trend som visar att databasstorlek
i sig förändrade förbättringens storlek i denna mätserie.

Samtidigt bör resultaten tolkas i ljuset av studiens begränsningar: mätningarna
genomfördes i en kontrollerad miljö med en testenhet, stabil Wi-Fi och
inaktiverad cachelagring. Slutsatsen är därför att metoden är praktiskt relevant
och lågtrösklig att införa i liknande applikationer, men att generalisering till
andra miljöer och applikationstyper bör göras med försiktighet.

\newpage

%-------------------------------------------------------------------------------
\bibliographystyle{unsrt}
\bibliography{base}

\vspace{12pt}

\rednote{ \textbf{Viktigt:} Detta har varit en mall för ditt gymnasiearbete med
fokus på X/Y/Z-strukturen. Se till att:
\begin{itemize}
    \item All malltext är borttagen innan inlämning
    \item Ta bort EXEMPEL-figurer som du inte använder (behåll max 3-4 figurer
    totalt)
    \item Din frågeställning följer formatet "Är Y lämplig för att förbättra X
    för produkttyp Z?"
    \item Du har minst 30 mätningar för statistisk validitet
    \item All kod och data är versionshanterad och reproducerbar
    \item Resultat och diskussion är tydligt separerade
\end{itemize}

Misslyckas du att följa denna struktur kan arbetet riskera att inte bli godkänt.
}

\clearpage
\appendix
\onecolumn

\section{Bilaga: Rådata från CSV}

Tabell \ref{tab:raw_csv_data} visar innehållet i \texttt{page\_load\_times\_utf8.csv}. Tabellen genereras automatiskt från CSV-filen, så när filen uppdateras uppdateras tabellen vid nästa kompilering.

\setlength{\LTleft}{0pt}
\setlength{\LTright}{0pt}

\begin{longtable}{p{0.30\textwidth}r r p{0.34\textwidth}}
\caption{Rådata importerad direkt från CSV-filen} \label{tab:raw_csv_data} \\
\toprule
Trace & Iteration & Load (ms) & Timestamp \\
\midrule
\endfirsthead

\toprule
Trace & Iteration & Load (ms) & Timestamp \\
\midrule
\endhead

\midrule
\multicolumn{4}{r}{Fortsätter på nästa sida} \\
\midrule
\endfoot

\bottomrule
\endlastfoot

\csvreader[
    filter expr={test{\ifnumgreater{\thecsvinputline}{1}}},
]{page_load_times_utf8.csv}{}{
    \nolinkurl{\csvcoli} & \csvcolii & \csvcoliii & \nolinkurl{\csvcoliv}\\
}
\end{longtable}

%%%%%%%%%%%%%%%%%%%%%%%%%%%%%%%%%%%%%%%%%%%%%%%%%%%%%%%%%%%%%%%%%%%%%%%%%%%%%%%%
\end{document}
%%%%%%%%%%%%%%%%%%%%%%%%%%%%%%%%%%%%%%%%%%%%%%%%%%%%%%%%%%%%%%%%%%%%%%%%%%%%%%%%

%%  LocalWords:  endnotes includegraphics fread ptr nobj noindent %  LocalWords:
%pdflatex acks
